
% ==============================================================================
% Master's Thesis LaTeX (single-file, no external \input required)
% Save as: thesis.tex
% Compile with: pdflatex thesis.tex  (run twice for references)
% ==============================================================================
\documentclass[12pt,a4paper,twoside,openright]{report}

% --- Page layout ---
\usepackage[a4paper,inner=35mm,outer=25mm,top=30mm,bottom=30mm]{geometry}

% --- Fonts and microtypography ---
\usepackage[T1]{fontenc}
\usepackage[utf8]{inputenc}
\usepackage{lmodern}
\usepackage{microtype}

% --- Math, tables, figures ---
\usepackage{amsmath,amssymb}
\usepackage{graphicx}
\usepackage{booktabs}
\usepackage{array}
\usepackage{tabularx}
\usepackage{float}
\usepackage{siunitx}
\usepackage{grffile} % improves robustness for file names with spaces

% --- Captions ---
\usepackage[font=small,labelfont=bf]{caption}

% --- Lists and spacing ---
\usepackage{enumitem}
\usepackage{setspace}

% --- Algorithms (pseudocode) ---
\usepackage[ruled,vlined]{algorithm2e}
\DontPrintSemicolon

% --- Headers/footers ---
\usepackage{fancyhdr}
\pagestyle{fancy}
\fancyhf{}
\fancyhead[LE,RO]{\thepage}
\fancyhead[LO]{\nouppercase{\leftmark}}
\fancyhead[RE]{\nouppercase{\rightmark}}
\renewcommand{\headrulewidth}{0.4pt}
\setlength{\headheight}{14.5pt}
\addtolength{\topmargin}{-2.5pt}

% --- Hyperlinks ---
\usepackage{hyperref}
\hypersetup{
  colorlinks=true,
  linkcolor=blue,
  citecolor=blue,
  urlcolor=blue,
  pdftitle={A Minimal-Input Framework for Cut-In Detection and Pair-Specific Risk Analysis in Highway Trajectory Data},
  pdfauthor={Sandeep Arsule and Shradha Shinde},
}

% --- SI units setup ---
\sisetup{
  detect-all,
  per-mode=symbol,
  range-phrase=--,
  range-units=single
}

% --- Utilities ---
\newcommand{\safeincludegraphics}[2][]{%
  \IfFileExists{\detokenize{#2}}{%
    \includegraphics[#1]{\detokenize{#2}}%
  }{%
    \fbox{%
      \parbox[c][0.25\textheight][c]{0.9\linewidth}{%
        \centering
        Figure file not found:\par
        \texttt{\detokenize{#2}}%
      }%
    }%
  }%
}

% --- Thesis metadata (edit these) ---
% \newcommand{\ThesisTitle}{Cut-in Risk Analysis from Highway Trajectory Data}
\newcommand{\ThesisTitle}{A Minimal-Input Framework for Cut-In Detection and Pair-Specific Risk Analysis in Highway Trajectory Data}
\newcommand{\FrameworkName}{Traj2Rel-SFC}
\newcommand{\ThesisSubtitle}{\FrameworkName: Trajectory-to-Relation Reconstruction and Context Signatures for Cut-In Interactions}
\newcommand{\AuthorOne}{Sandeep Arsule}
\newcommand{\AuthorTwo}{Shradha Shinde}
\newcommand{\SupervisorName}{Junyi Gu, Ph.D.}
\newcommand{\ExaminerName}{Aris Alissandrakis}
\newcommand{\UniversityName}{Chalmers University of Technology}
\newcommand{\DepartmentName}{Department of Computer Science and Engineering}
\newcommand{\CityCountry}{Gothenburg, Sweden}
\newcommand{\ThesisDate}{\today}

% Optional: link to your reproducibility package / repository
\newcommand{\ArtifactURL}{\url{https://github.com/arsulesandy/cutin-risk-analysis}}

% ==============================================================================
\begin{document}

% ==============================================================================
% Front matter
% ==============================================================================
\pagenumbering{roman}

% --- Title page ---
\begin{titlepage}
\thispagestyle{empty}
\vspace*{2cm}

\begin{center}
{\Large \textsc{\UniversityName}}\\[0.5cm]
{\large \DepartmentName}\\[2.0cm]

{\huge \bfseries \ThesisTitle}\\[0.6cm]
{\Large \ThesisSubtitle}\\[2.0cm]

{\large Master's Thesis in Software Engineering (30 credits)}\\[1.2cm]

{\large
\textbf{Author 1:} \AuthorOne\\
\textbf{Author 2:} \AuthorTwo
}\\[1.2cm]

{\large
\textbf{Supervisor:} \SupervisorName\\
\textbf{Examiner:} \ExaminerName
}\\[1.6cm]

{\large \CityCountry}\\
{\large \ThesisDate}
\end{center}

\vfill
\end{titlepage}

% --- Abstract ---
\cleardoublepage
\chapter*{Abstract}
\addcontentsline{toc}{chapter}{Abstract}

Lane ids and neighbour ids are important when studying vehicle interactions in trajectory datasets. The problem is that many analysis pipelines depend on these derived fields directly, which makes them hard to reuse when the fields are unavailable, inconsistent, or defined differently across datasets. This thesis presents \FrameworkName{}, a minimal-input framework that reconstructs lane assignment and same-lane predecessor/follower relations from trajectory geometry ($x,y$) and lane-marking metadata.

Based on these reconstructed relations, the framework detects cut-ins as explicit cutter--follower pairs, computes distance headway (DHW), time headway (THW), and time-to-collision (TTC) for that pair, and represents the surrounding traffic context with an ego-centric occupancy grid. To make the context comparable across travel directions, the grid is mirror-normalised and encoded as a reversible 16-bit Hilbert space-filling-curve signature.

The framework is evaluated on the \textit{highD} dataset, where highD lane and neighbour identifiers are used only as reference labels and not as method inputs. Across recordings 01--60, same-lane predecessor/follower reconstruction with inferred lanes reaches mean accuracy 0.999810/0.999811, and cut-in events detected from reconstructed relations match the reference-label events with $F_1 = 0.998938$. The detector extracts 10,026 cut-ins, of which 10,023 remain after stage-window validity checks. For 1,421,980 stage rows, reconstructed and reference context signatures agree exactly in 99.49\% of cases, and the SFC encoding is additionally verified through exhaustive round-trip testing over all $2^{16}$ code values.

The main contribution of the thesis is not only that the method matches highD closely, but that it keeps the interacting pair, the indicator computation, and the local context representation explicit. This makes the analysis reproducible, easier to inspect, and reusable on data sources that do not provide neighbour relations directly.

\noindent\textbf{Keywords:} cut-in detection, minimal-input reconstruction, pair-specific risk analysis, surrogate safety measures, interaction context signature, space-filling curve, highD, reproducibility.

% --- Acknowledgements ---
\cleardoublepage
\chapter*{Acknowledgements}
\addcontentsline{toc}{chapter}{Acknowledgements}

We thank our supervisor for continuous feedback and for pushing us to define a clear system boundary: what information the framework is allowed to use as input, and what information must be treated as evaluation labels. We also thank the course staff and seminar group for discussions about scope, evaluation, and writing. Finally, we are grateful to the authors of the highD dataset for making the dataset available to the research community.

% --- Contents / lists ---
\cleardoublepage
\tableofcontents

\cleardoublepage
\addcontentsline{toc}{chapter}{\listfigurename}
\listoffigures

\cleardoublepage
\addcontentsline{toc}{chapter}{\listtablename}
\listoftables

% --- Abbreviations ---
\cleardoublepage
\chapter*{Abbreviations}
\addcontentsline{toc}{chapter}{Abbreviations}

\begin{tabular}{@{}ll@{}}
\toprule
Abbreviation & Meaning \\
\midrule
DHW & Distance headway (longitudinal gap) \\
THW & Time headway (gap divided by follower speed) \\
TTC & Time-to-collision (gap divided by closing speed, if closing) \\
SSM & Surrogate safety measure \\
SFC & Space-filling curve \\
XY-lane & Lane inference from $y$ and lane markings \\
\bottomrule
\end{tabular}

% ==============================================================================
% Main matter
% ==============================================================================
\cleardoublepage
\setcounter{page}{1}
\pagenumbering{arabic}
\onehalfspacing

% ==============================================================================
\chapter{Introduction}
\label{chap:introduction}

Lane changes are common on highways and are not necessarily unsafe. A \emph{cut-in} is a lane change where a vehicle moves into a lane directly in front of another vehicle, thereby creating a direct longitudinal interaction in the target lane. Depending on available gap, relative speed, and driver behaviour, a cut-in can lead to uncomfortable or safety-critical situations.

Trajectory datasets enable large-scale analysis of such interactions. However, many datasets provide surrogate safety measures (e.g., TTC/THW/DHW) only for a \emph{predefined} neighbour relationship (typically ``ego--same-lane predecessor''). If we analyse cut-ins but reuse indicators computed for a different pair, the analysis becomes internally inconsistent even if computations are correct.

\section{System boundary: analysing a dataset vs evaluating a method}
A central writing and evaluation decision in this thesis is the \textbf{system boundary}. With a dataset such as \textit{highD}, there are two valid perspectives:

\begin{itemize}[leftmargin=1.2em]
  \item \textbf{Dataset analysis perspective:} treat highD's derived fields (lane ids, neighbour ids, and possibly TTC/THW/DHW) as part of the input, and use them to analyse traffic behaviour in highD.
  \item \textbf{Method-evaluation perspective:} treat highD's derived fields as \emph{reference labels} , and evaluate a method that must work using only lower-level inputs such as trajectories and lane markings.
\end{itemize}

This thesis follows the \textbf{method-evaluation perspective}. Concretely, our \emph{system under evaluation} reconstructs lane indices and neighbour relations from $x,y$ trajectories and lane markings only. HighD's lane and neighbour identifiers are used \emph{only} as reference labels for evaluation and sanity checking.

\section{Problem statement}
This thesis addresses the lack of a reproducible framework that can rebuild lane and same-lane relations from minimal inputs, define cut-ins as explicit cutter--follower events, compute safety indicators for that pair, and represent the surrounding traffic scene in a compact form.

\section{Why matching highD is not enough}
High agreement against highD is important, but by itself it is not a full thesis contribution. If the method only reproduces the dataset's provided relations, then the work can still look like a careful reconstruction exercise. The thesis therefore also asks what becomes possible once the relevant pair is reconstructed explicitly. In our case, that means defining a cut-in as a cutter--follower event, computing DHW, THW, and TTC for that pair, and showing how the interpretation changes across stages and thresholds. This is the part that turns reconstruction accuracy into a useful analysis pipeline rather than only a matching exercise.

\section{Key idea and novelty}
The main idea behind \FrameworkName{} is simple: first reconstruct lane and same-lane neighbour relations from minimal inputs, then use those reconstructed relations to define cut-ins in a clear way, and finally encode local context so that events can be compared later. The framework therefore has three connected parts:

\begin{enumerate}[leftmargin=1.6em]
  \item a \textbf{trajectory-to-relation reconstruction layer} (Traj2Rel) that rebuilds lane assignment and same-lane neighbour relations from minimal geometric inputs ($x,y$ and lane markings),
  \item an \textbf{explicit interaction-pair layer} that defines a cut-in as a cutter--follower event with a reproducible onset and pair-consistent indicator computation, and
  \item a \textbf{context-signature layer} that turns the surrounding traffic configuration into a compact, reversible integer code using a Hilbert space-filling-curve traversal.
\end{enumerate}

The key design choice is the evaluation boundary. highD lane and neighbour identifiers are used only as reference labels and never as hidden inputs to the method. This lets us test whether the minimal-input path works on its own and keeps the later cut-in analysis independent from dataset-specific relation columns. The individual ingredients are known in the literature; the contribution here is to bring them together in one clear pipeline for cut-in analysis and to evaluate that pipeline explicitly against reference labels.

\section{Goal}
The goal is to build and evaluate a reproducible framework that:
\begin{enumerate}[leftmargin=1.6em]
  \item infers lane ids from $y$ positions and lane markings (XY-lane),
  \item reconstructs same-lane preceding/following relations from geometry ($x$ ordering),
  \item detects lane changes and cut-ins and identifies the affected follower for each cut-in,
  \item computes DHW/THW/TTC for the correct cutter--follower pair,
  \item encodes stage-wise spatial context around cut-ins in a direction-consistent way using SFC representations,
  \item evaluates reconstruction quality and event consistency against highD reference labels, and
  \item makes the cut-in definition explicit through stage definitions, threshold sensitivity, and indicator-disagreement analysis.
\end{enumerate}

\section{Research questions}
\textbf{RQ1:} How accurately can lane indices and same-lane neighbour relations be reconstructed from trajectories and lane markings, compared to highD's provided lane and neighbour identifiers?

\textbf{RQ2:} How can lane changes and cut-ins be detected using only reconstructed relations, how consistent are the detected cut-in events with a reference-label path that uses highD relations, and what does explicit cutter--follower pairing add to risk-feature construction?

\textbf{RQ3:} How can spatial context around cut-ins be encoded using SFC signatures, and how can the resulting representation support clear, direction-consistent context comparison against highD's reference relations?

\section{Scope and delimitations}
We focus on the highD dataset (highway traffic, drone-based trajectories) and offline analysis. The thesis does not aim to infer driver intent beyond observable trajectories, does not include simulation-based re-enactment, and does not claim that any single risk threshold is universally valid. Instead, the emphasis is on scenario extraction, metric correctness, validation against reference labels when applicable, and reusable feature generation.

\section{Relation to the accepted proposal}
This thesis is based on an accepted proposal for a ``Framework for Cut-In Vehicle Detection and Lane-Change Risk Analysis in Naturalistic Highway Trajectories''. During implementation, the scope was refined to emphasise (i) \emph{minimal-input reconstruction} of lanes and neighbour relations, (ii) \emph{explicit cutter--follower pairing} and pair-consistent indicator computation, (iii) \emph{direction-normalised} spatial context encoding and verification, and (iv) a \emph{clear separation} between event extraction and downstream risk labelling. The accepted proposal also mentioned computational efficiency, but the final thesis makes a stronger and better-supported contribution by focusing on method reuse, reference-label validation, and compact context encoding rather than on runtime benchmarking alone. The revised research questions (RQ1--RQ3) remain aligned with the proposal's core aim: reproducible detection and interpretable risk analysis at scale \cite{highd_2018,wang2021ssm}.

\section{Contributions}
This thesis presents \FrameworkName{}, a reconstruction-and-signature framework for cut-in interactions, and makes the following contributions:

\begin{enumerate}[leftmargin=1.6em]
  \item \textbf{Minimal-input relation reconstruction:} lane inference (XY-lane) and same-lane predecessor/follower reconstruction from $x,y$ trajectories and lane markings, evaluated frame by frame against highD reference labels.
  \item \textbf{Explicit cutter--follower cut-in definition:} cut-ins are detected as explicit interacting pairs with a reproducible onset definition ($t_0$), which makes later feature extraction consistent and easy to inspect.
  \item \textbf{Pair-specific indicator analysis:} DHW, THW, and TTC are computed for the actual cutter--follower pair and interpreted stage-wise, using literature-backed reasoning for why TTC and THW are useful in different parts of the interaction.
  \item \textbf{Reversible context signatures:} a direction-normalised semantic occupancy grid is encoded into a reversible 16-bit Hilbert SFC signature and verified against reference occupancy built from highD neighbour relations.
  \item \textbf{Clear evaluation setup and reproducible outputs:} a two-path design separates the minimal-input reconstruction path from the reference-label path, and the full run produces deterministic tables and figures that can be regenerated from the same settings.
\end{enumerate}

\section{Thesis outline}

Chapter~\ref{chap:background} introduces the dataset and core concepts. Chapter~\ref{chap:relatedwork} presents the related work. Chapter~\ref{chap:researchmethod} describes the research method and evaluation approach. Chapter~\ref{chap:theory} defines events, stages, indicators, and SFC signatures. Chapter~\ref{chap:methods} details the framework and validation. Chapter~\ref{chap:results} reports results. Chapter~\ref{chap:discussion} discusses limitations and implications, and Chapter~\ref{chap:conclusion} concludes.

% ==============================================================================
\chapter{Background}
\label{chap:background}

\section{Trajectory data and the highD dataset}
We use the \textit{highD} dataset, a drone-based collection of naturalistic highway trajectories \cite{highd_2018}. HighD provides per-frame vehicle states (position, velocity, acceleration), lane identifiers, and neighbour identifiers such as same-lane preceding and following vehicles. The recordings are sampled at \SI{25}{\hertz}.

A recording is distributed as three data tables: per-frame trajectories, per-vehicle metadata, and recording-level metadata. For completeness, the corresponding file names in the released dataset are listed in Appendix~\ref{app:field-mapping}. Throughout the thesis we refer to these sources by their role (trajectories / track metadata / recording metadata) to keep the presentation readable.



\section{Neighbour relations in highD}
In addition to trajectories, highD provides identifiers of surrounding vehicles for each vehicle and time step (e.g., the closest vehicle ahead and behind in the same lane, and corresponding neighbours in adjacent lanes). These identifiers are convenient for analysing highD itself and, in this thesis, serve as \emph{reference labels} for evaluating the reconstruction method. The reconstruction path described in this thesis does not use these identifiers as inputs; it derives neighbour relations from geometry and compares them against the highD-provided identifiers. A mapping between the terminology used in this thesis and the relevant highD columns is provided in Appendix~\ref{app:field-mapping}.

\section{Why compute indicators ourselves}
For cut-in risk, the key interaction is between the cutter and the follower that ends up behind the cutter in the target lane. Even if the dataset provides surrogate measures, we compute DHW/THW/TTC ourselves for the cutter--follower pair to ensure that analysis reflects the intended interaction and to support framework variants where neighbours and lanes are reconstructed from $x,y$ and lane markings.

% ==============================================================================
\chapter{Related Work}
\label{chap:relatedwork}

This chapter reviews (i) surrogate safety measures, (ii) lane-changing and cut-in analysis, (iii) highD and dataset constraints, and (iv) compact context encodings for interaction analysis.

\section{Literature identification approach}
This thesis is not presented as a full systematic literature review. However, to keep the related-work selection transparent and reproducible, we followed a structured search-and-selection approach inspired by evidence-based guidelines \cite{kitchenham2007,wohlin2014}. We searched IEEE Xplore, the ACM Digital Library, and Google Scholar (with Scopus used when available) for peer-reviewed work on (i) surrogate safety measures, (ii) lane changes and cut-ins, (iii) \textit{highD} and similar trajectory datasets, and (iv) spatial/context representations for traffic interactions. The search emphasised recent work relevant to \textit{highD} (primarily 2018--2025) while also including seminal papers needed to define key concepts.

Example search strings used during the search include:
\begin{quote}\small\ttfamily\raggedright
("cut-in" OR "lane change") AND (trajectory OR dataset) AND\\
(risk OR safety)\\[2pt]
"highD" AND ("lane change" OR "cut-in") AND\\
(TTC OR THW OR "surrogate safety")\\[2pt]
("surrogate safety measure" OR "traffic conflict") AND\\
(validation OR review)\\[2pt]
("occupancy grid" OR "spatial context") AND\\
trajectory AND vehicle\\[2pt]
("space-filling curve" OR Hilbert OR Morton) AND\\
(occupancy OR grid OR encoding)
\end{quote}

We included papers that address highway lane changes/cut-ins, define or review surrogate safety measures, use \textit{highD} (or closely related datasets) for scenario extraction, or propose compact spatial context representations for interaction analysis. We complemented keyword searching with backward and forward snowballing on core papers (e.g., the \textit{highD} dataset paper and SSM review papers) \cite{wohlin2014}.

\section{Surrogate safety measures and validation}
Surrogate safety measures (SSMs) are used because crashes (collisions) are rare in naturalistic data. TTC was introduced early as a near-miss concept \cite{hayward1972}. SSAM operationalised a set of commonly used indicators for conflict analysis and provided tooling for systematic application \cite{ssam2008}. Extended TTC formulations have been proposed to adapt TTC-style reasoning to broader traffic conditions \cite{minderhoud2001}. Recent reviews and validation-oriented studies emphasise that SSM values depend strongly on modelling assumptions, measurement noise, and traffic state, and therefore need careful interpretation \cite{wang2021ssm,johnsson2021validation}.

\section{Lane-changing and cut-in analysis}
Lane changes and cut-ins are often described through a combination of incentive, safety, and temporal structure. Classical microscopic models such as the Intelligent Driver Model (IDM) formalise desired spacing and speed adaptation in car-following \cite{treiber2000idm}. Lane-changing models such as MOBIL make safety and incentive criteria explicit \cite{kesting2007mobil}.

Naturalistic cut-in studies show that cut-ins are not just lane changes with a binary flag. Wang et al. develop an extraction algorithm for 5,608 cut-in events and report substantial variation in cut-in duration (0.7--12.4 s) together with gap-acceptance effects from environment, vehicle type, and kinematics \cite{wang2019cutin}. Zhou and Itoh study driver risk perception in two lane-change stages and find that drivers tend to rely on TTC when deciding to change lanes and on both TTC and THW while cutting into the target lane \cite{zhou2016riskperception}. This directly motivates phase-aware interpretation of surrogate measures in the present thesis.

Recent risk-prediction frameworks further reinforce the need for explicit scenario construction. On highD, Shangguan et al. integrate lane-change intention recognition with risk prediction and show that interaction characteristics with surrounding vehicles are central to risk estimation \cite{shangguan2022}. Liu et al. jointly predict lane-change manoeuvre, time-to-lane-change, and risk level, again using highD as the evaluation basis \cite{liu2025multitask}. These studies are predictive rather than reconstructive, but they support the thesis decision to preserve explicit pair identities, stage alignment, and interaction context features.

For cut-in and lane-merging behaviour in automated driving, structured phase models with explicit risk logic are also common. For example, Hwang et al. propose a finite-state-machine approach for lane-merging cut-ins with TTC-based risk handling \cite{hwang2022fsmrl}. While this thesis does not implement a controller, such phase-based framing motivates clear scenario segmentation and transparent risk definitions for offline analysis.

\section{HighD and dataset constraints}
HighD provides high-precision trajectories and neighbour relations for large-scale scenario extraction \cite{highd_2018}. At the same time, dataset conventions and lane assignment noise can affect how lane-change timing is interpreted. Recent work discusses lane-change inconsistencies in highD and motivates careful validation when reconstructing events from lane ids and geometry \cite{wanggao2025highd}.

\section{Context encoding for interaction analysis}
Scalar indicators (e.g., minimum THW) compress interactions to a few values and may miss relevant local traffic structure. Grid-based representations such as occupancy grids are widely used for compactly describing local surroundings \cite{elfes1989}. For trajectory prediction and interaction modelling, pooling neighbour information into a spatial grid can improve robustness and interpretability, including for vehicle settings \cite{deo2018}. 

To compress spatial grids into a compact signature while preserving locality, space-filling curves (SFCs) such as the Hilbert curve provide a principled ordering of 2D cells into 1D sequences \cite{hilbert1891,moon2001}. This thesis uses a small ego-centric 3$\times$3 occupancy grid and encodes it into a 16-bit SFC signature to obtain a reversible integer-code representation of local context.

\section{Summary and research gap}
Prior work establishes multiple surrogate safety measures, phase models, and predictive lane-change frameworks, but three gaps motivate this thesis. First, cut-in analysis requires explicit pairing between the cutter and the impacted follower; dataset-provided SSMs may not correspond to that pair. Second, existing highD-based risk models usually assume that scenario samples and interaction partners already exist; they do not solve the minimal-input reconstruction problem when lane and neighbour identifiers are unavailable. Third, compact context representations are rarely combined with end-to-end reference-label verification. This thesis addresses these gaps via a reproducible framework that reconstructs lane and neighbour relations from minimal inputs, computes pair-consistent and phase-aware indicators, and produces verified SFC-based context signatures.

% ==============================================================================
\chapter{Research Method}
\label{chap:researchmethod}

\section{Methodological framing}
This thesis follows a design-oriented empirical methodology. We propose a reconstruction and representation method that derives lane indices and same-lane predecessor/follower relations directly from geometric trajectory information. The \textit{highD} dataset is treated as a reference source for evaluation: highD-provided lane and neighbour identifiers are used as reference labels, not as inputs to the method under evaluation.

\section{Evaluation logic: correctness and usefulness}
The evaluation in this thesis has two jobs. First, it must show that the minimal-input reconstruction is correct enough to reproduce highD's relational structure. Second, it must show why the framework is useful once that correctness baseline is established. For that reason, the evaluation is reported at four levels:

\paragraph{Frame-level validation (relations).}
For every frame where a relation exists, the predicted lane assignment and same-lane leader/follower identifiers are compared against the corresponding highD reference labels.

\paragraph{Event-level validation (cut-ins).}
Cut-in events detected from reconstructed relations are compared against cut-ins detected using reference relations. Agreement is reported using precision, recall, and $F_1$ score with an onset tolerance in frames.

\paragraph{Representation-level validation (occupancy and SFC signatures).}
For cut-in stage frames, occupancy grids are constructed both from reconstructed relations and from reference relations. The resulting binary grids and their derived SFC codes are compared using cell-wise agreement and Hamming distance. This checks whether the context signatures preserve the same neighbour-occupancy structure as the reference path.

\paragraph{Usefulness beyond agreement (pairing, stages, and indicators).}
Beyond reference-label agreement, we evaluate whether the framework leads to a clear and interpretable cut-in analysis once the interacting pair has been reconstructed. This part is reported through stage-based summaries, risk-threshold sensitivity, and TTC-versus-THW disagreement analysis. The goal is not to claim one universal threshold, but to show why explicit pair and stage definitions matter.

\paragraph{Qualitative verification (interactive dashboard).}
To complement the quantitative metrics, we performed qualitative spot checks using an interactive dashboard that synchronises video playback with reconstructed relations, event timing, and context signatures. This helped verify that the reported metrics correspond to visually sensible interactions and supported illustrative examples.

\section{Metrics}
Lane and neighbour reconstruction are evaluated using exact-match accuracy. Because a large fraction of frames contain no neighbour (encoded as id $0$), we report accuracy both over all frames and conditional on frames where the reference id is non-zero; the latter avoids inflated scores from trivial no-neighbour matches. Cut-in detection agreement is evaluated using precision, recall, and $F_1$ score. Occupancy agreement is evaluated using:

\begin{itemize}[leftmargin=1.2em]
  \item \textbf{Cell-wise match ratio} for the 3$\times$3 semantic grid:
  \[
    \text{Match Ratio}=\frac{1}{9}\sum_{i=1}^{9}\mathbf{1}\big(G_r[i]=G_{ref}[i]\big),
  \]
  where $G_r$ is the reconstructed grid and $G_{ref}$ is the reference grid.
  \item \textbf{Hamming distance} between the 16-bit SFC codes (bitwise differences).
  \item \textbf{Perfect-match rate}, i.e., percentage of frames with identical grids/codes.
\end{itemize}

For the operational-usefulness layer, we additionally report event counts after acceptance criteria, risk-prevalence sensitivity under alternative THW thresholds, and TTC coverage (finite versus infinite) to show how explicit pair and stage definitions change interpretation.

\section{Ethical and legal considerations}
The thesis uses an existing trajectory dataset collected from public roads. No new data is collected and no human subjects are recruited. Derived outputs should be shared only in ways that respect the dataset licence and terms of use.

\section{Use of Generative AI Tools}
Generative AI tools (e.g., large language models) were used in a limited and supportive role during this project. AI assistance was used primarily to troubleshoot development environment configuration issues, identify and resolve Python runtime errors and dependency conflicts, and suggest refactoring patterns to standardise code structure and improve readability. AI assistance was also used to draft diagram code (e.g., Mermaid), which was then manually reviewed and exported as static figures for inclusion in the thesis.

All core research ideas, problem formulations, definitions (e.g., cut-in detection logic, stage modelling, and pair-consistent metric computation), methodological decisions, evaluation design, and interpretation of results were developed independently by the authors. AI-generated suggestions were treated as technical assistance similar to consulting documentation or community forums, and any suggested code changes were manually reviewed, validated, and tested before integration. The authors take full responsibility for the correctness of the implementation, the integrity of the analysis, and the conclusions presented in this thesis.

\chapter{Theory}
\label{chap:theory}

\section{Lane change vs cut-in}
\subsection{Lane change}
We define a \textbf{lane change event} as a transition from one stable lane block to another stable lane block. ``Stable'' means the vehicle remains in that lane for at least a minimum number of consecutive frames before and after the transition.

\subsection{Cut-in}
We define a \textbf{cut-in event} as a lane change where the lane-changing vehicle (the \emph{cutter}) becomes the same-lane predecessor of another vehicle (the \emph{follower}) in the target lane. This definition yields an explicit interacting pair $(A,B)$ with cutter $A$ and follower $B$.

\section{Cut-in onset and stage model}
We define cut-in onset $t_0$ as the first frame where the follower's same-lane predecessor becomes the cutter id (i.e., the first frame where the follower relationship is established). Time around $t_0$ is segmented into four stages:
\begin{itemize}[leftmargin=1.2em]
  \item \textbf{Intention:} $[t_0-\SI{4}{\second},\, t_0-\SI{2}{\second})$
  \item \textbf{Decision:} $[t_0-\SI{2}{\second},\, t_0)$
  \item \textbf{Execution:} $[t_0,\, t_0+\SI{2}{\second})$
  \item \textbf{Recovery:} $[t_0+\SI{2}{\second},\, t_0+\SI{4}{\second}]$
\end{itemize}

\begin{figure}[H]
\centering
\safeincludegraphics[width=\textwidth]{fig_stage_timeline.pdf}
\caption{Stage model around the cut-in onset $t_0$ used for stage-wise feature extraction and comparison.}
\label{fig:stage-timeline}
\end{figure}

The four-stage decomposition is an operational definition for offline analysis rather than a claim about exact human cognitive phases. It is nevertheless consistent with structured lane-change and cut-in formulations in the literature, where intention, decision, insertion, and recovery are treated as distinct but related parts of the manoeuvre \cite{hwang2022fsmrl,zhou2016riskperception}.

\section{Surrogate safety measures}
Let $A$ be the leader (cutter) and $B$ the follower. We use:
\begin{itemize}[leftmargin=1.2em]
  \item \textbf{DHW} (distance headway): longitudinal gap between leader rear and follower front.
  \item \textbf{THW} (time headway): $\mathrm{THW} = \mathrm{DHW} / v_B$ (with safe handling of small $v_B$).
  \item \textbf{TTC} (time-to-collision): under constant-velocity assumptions,
  \[
    \mathrm{TTC}=
    \begin{cases}
      \mathrm{DHW}/(v_B - v_A), & \text{if } v_B > v_A \\
      \infty, & \text{otherwise.}
    \end{cases}
  \]
\end{itemize}

Because cut-ins change the relevant interacting pair, all three measures are computed for the explicit cutter--follower pair rather than reusing dataset-provided values for a possibly different neighbour pair. Recent reviews emphasise that surrogate safety measures are definition-dependent and should be interpreted with care \cite{wang2021ssm,johnsson2021validation}. In addition, Zhou and Itoh report that TTC is particularly salient to drivers when deciding whether to merge, while TTC and THW together remain informative during the actual insertion phase \cite{zhou2016riskperception}. We therefore compute DHW, THW, and TTC throughout the stage-aligned window and interpret them as complementary rather than interchangeable indicators.

\section{Risk label}
Because risk is not standardised, we use a transparent label rule for descriptive comparison:
\[
  \text{risky} := \min(\mathrm{THW}) \text{ in execution stage} < \SI{0.7}{\second}.
\]
This rule is not presented as a universal definition of risk. We choose execution-stage THW for this descriptive label because it remains defined after the new follower relation has been established, even when TTC becomes infinite due to non-closing conditions. The thesis therefore reports threshold sensitivity to keep the operational choice explicit rather than implicit \cite{wang2021ssm,zhou2016riskperception}.

\section{Interaction context signatures (SFC encoding)}
For each frame around a cut-in event, we define an ego-centric \textbf{3$\times$3 semantic occupancy grid} around the cutter:
\begin{itemize}[leftmargin=1.2em]
  \item columns: \{left lane, same lane, right lane\},
  \item rows: \{preceding, alongside, following\}.
\end{itemize}

Each cell is treated as binary occupancy ($1$ if at least one vehicle occupies the cell, else $0$). The 3$\times$3 grid is embedded in a fixed 4$\times$4 grid and encoded as a 16-bit integer using a locality-preserving space-filling curve order (Hilbert order by default). Because the mapping is deterministic and reversible, we can decode a code back into a grid.

\begin{figure}[H]
\centering
\safeincludegraphics[width=\textwidth]{fig_sfc_signature_overview.pdf}
\caption{Interaction context signatures. A 3$\times$3 ego-centric semantic occupancy grid is constructed around the cutter, mirror-normalised to a canonical direction, embedded into a fixed 4$\times$4 grid, and encoded into a reversible 16-bit space-filling-curve (Hilbert) code.}
\label{fig:sfc-signature-overview}
\end{figure}

To verify that the encoding itself is correct, we perform an exhaustive round-trip test across all $2^{16}$ possible bit patterns, confirming that encoding followed by decoding returns the original grid for both Hilbert and Morton (Z-order) traversals.

% ==============================================================================
\chapter{Methods}
\label{chap:methods}

\section{Framework overview}
\FrameworkName{} is a modular framework for \emph{trajectory-to-relation reconstruction} and \emph{interaction-context signatures}. The core computation uses only geometric trajectory signals ($x,y$, velocities, vehicle dimensions) and lane-marking metadata. HighD-provided lane and neighbour identifiers are treated as \textbf{reference labels} and are used only inside evaluation components.

\subsection*{Two-path execution for reference-label evaluation}
A central design choice is to execute the same analysis logic in two parallel paths:

\begin{itemize}[leftmargin=1.2em]
  \item \textbf{Reference-label path:} uses highD-provided lane ids and neighbour ids to define reference relations, reference cut-in events, and reference occupancy signatures.
  \item \textbf{Reconstruction path (system under evaluation):} uses only $x,y$ trajectories and lane markings to infer lane indices (XY-lane), reconstruct same-lane neighbours, and then produce the same event and signature outputs \emph{without} using highD neighbour columns as inputs.
\end{itemize}

This design treats highD as a reference source for evaluation rather than a dependency of the method. It also enables direct, like-for-like agreement measurements at three levels: \textbf{frame-level} (lane and neighbour relations), \textbf{event-level} (cut-in detection), and \textbf{representation-level} (occupancy matrices and SFC signatures).

\subsection*{Macro-phases and component naming}
To keep the method presentation readable (and to avoid numeric ``01/02/03'' style narration), we group the framework into macro-phases with descriptive names. The reconstruction path can be interpreted as a \emph{sequence of deterministic processing stages}: each component applies explicit acceptance criteria (e.g., lane stability, relation duration, stage-window validity) and emits traceable outputs.

Figure~\ref{fig:framework-overview} illustrates the framework and the reconstruction path.


\begin{figure}[H]
    \centering
    \safeincludegraphics[width=1.1\linewidth]{main_figure.png}
\caption{Overview of \FrameworkName{} and its reference-label evaluation design. The reconstruction path infers lane membership and same-lane relations from trajectories and lane markings; highD lane and neighbour identifiers are used only as reference labels for agreement measurement.}
    \label{fig:framework-overview}
\end{figure}


\begin{table}[H]
\centering
\caption{Macro-phases and outputs of \FrameworkName{}.}
\label{tab:framework-phases}
{\small
\setlength{\tabcolsep}{4pt}
\begin{tabularx}{\textwidth}{@{}l l >{\raggedright\arraybackslash}X@{}}
\toprule
Phase & Component & Main outputs \\
\midrule
Preliminary processing & Data harmonisation & Consolidated per-frame dataset; sanity checks. \\
\midrule
Core reconstruction & XY-lane inference (Traj2Rel) & Predicted lane index $\hat{L}$ per frame from $y$ and lane markings. \\
Core reconstruction & Neighbour reconstruction (Traj2Rel) & Predicted same-lane leader and follower identifiers per frame via within-lane longitudinal ordering. \\
\midrule
Interaction mining & Cut-in event extraction & Cut-in table with explicit cutter--follower pairs and onset $t_0$. \\
Safety quantification & Pair-consistent indicators \& stages & DHW/THW/TTC for the interacting pair; stage-wise aggregates and risk summaries. \\
\midrule
Context signatures & Occupancy \& SFC encoding & Mirror-normalised occupancy matrices; reversible 16-bit SFC context signatures (hex + bits); stage-level context features. \\
\midrule
Reference-label evaluation & Agreement reporting & Frame-level accuracy, event-level precision/recall/$F_1$, and representation-level match/Hamming-distance diagnostics. \\
\bottomrule
\end{tabularx}}
\end{table}

The remainder of this chapter describes each component and the corresponding reference-based validation checks.


\subsection{Qualitative inspection support}
Quantitative agreement metrics provide an overall assessment of reconstruction quality, but they do not on their own reveal whether failures are systematic (e.g., due to direction conventions, onset misalignment, or lane-boundary interpretation). To support qualitative verification, we developed an interactive inspection dashboard that synchronises (i) aerial video playback and track overlays, (ii) the currently selected cut-in event and its cutter--follower pair, and (iii) the derived stage window, indicators, and context signature for the current frame.

The dashboard was used to inspect representative cases and extremes (e.g., lowest-TTC events) and to verify that reconstructed relations and occupancy signatures correspond to the visible scene. It is a support tool for verification and communication rather than a standalone thesis contribution: it was not used to manually correct events or labels, and all reported results are produced by the framework's deterministic components.


\begin{figure}[H]
  \centering
  \safeincludegraphics[width=\textwidth]{fig_gui_command_deck.png}
  \caption{Interactive inspection dashboard used for qualitative verification and illustration. The interface links video playback with reconstructed lane and neighbour relations, highlights the cutter--follower pair for selected cut-in events, and displays the corresponding stage and context signature for the current frame.}
  \label{fig:gui-dashboard}
\end{figure}


\section{Key configuration values}
Table~\ref{tab:config} summarises key configuration values used in the final run (highD 01--60, \SI{25}{\hertz}).

These values act as explicit cut-offs (e.g., stability requirements, search windows, and stage alignment) and are reported to support reproducibility and sensitivity analysis.

\begin{table}[H]
\centering
\caption{Key configuration values for the final framework run.}
\label{tab:config}
{\small
\setlength{\tabcolsep}{4pt}
\begin{tabularx}{\textwidth}{@{}l l X@{}}
\toprule
Component & Parameter & Value \\
\midrule
Lane-change detection & Min stable before/after & 25 frames ($\approx \SI{1.0}{\second}$) \\
Lane-change detection & Ignored lane ids & \{0\} \\
Cut-in detection & Search window & 50 frames ($\approx \SI{2.0}{\second}$) \\
Cut-in detection & Min relation duration & 10 frames ($\approx \SI{0.4}{\second}$) \\
Cut-in detection & Max relation delay & 15 frames ($\approx \SI{0.6}{\second}$) \\
Indicator computation & Min speed for THW & \SI{0.1}{\meter\per\second} \\
Indicator computation & Closing speed epsilon & $10^{-6}$ \\
Indicator computation & Window (pre/post) for diagnostics & 50/75 frames ($\approx \SI{2}{\second}/\SI{3}{\second}$) \\
Stage model & Total analysis window & $\pm \SI{4}{\second}$ around $t_0$ \\
Risk label & Execution THW threshold & \SI{0.7}{\second} \\
SFC encoding & Ahead/behind range & Disabled (no fixed cutoff) \\
SFC encoding & Alongside threshold & \SI{4.0}{\meter} (final Step 14 thesis run) \\
SFC encoding & Order & Hilbert \\
\bottomrule
\end{tabularx}}
\end{table}

These values are operational rather than universal. Stable blocks and pair-persistence criteria make the cut-in definition explicit, while the phase windows and indicator interpretation are consistent with literature that treats lane changing as a staged process and emphasises the context-dependence of gap acceptance and risk indicators \cite{wang2019cutin,zhou2016riskperception,hwang2022fsmrl}.


\section{Preliminary processing: data harmonisation}
Each highD recording is distributed as three data tables: (i) per-frame trajectories, (ii) per-vehicle metadata, and (iii) recording-level metadata. In preliminary processing we load these tables, perform basic consistency checks (e.g., frame rate and track bounds), and combine them into a consolidated per-frame dataset where each row corresponds to one vehicle at one time step together with the metadata needed in later phases. This consolidated representation is used by both the reference-label path and the reconstruction path so that subsequent comparisons are like-for-like.

We also extract lane-marking positions from the recording metadata because they provide the only map-related input used by the reconstruction method. Dataset-provided lane and neighbour identifiers are retained only for evaluation and are never used as inputs in the reconstruction path.

\section{Core reconstruction: lane and neighbour inference (Traj2Rel)}
The reconstruction layer (Traj2Rel) derives two pieces of relational information from minimal inputs: (i) lane membership for each vehicle and time step, and (ii) the identity of the closest vehicle ahead and behind in the same lane. Together these relations are sufficient to define leader--follower interactions and to express cut-in events as explicit cutter--follower pairs.

\subsection{Lane inference from lateral position (XY-lane)}
Lane membership is inferred by mapping lateral position to the interval between two consecutive lane markings provided in the recording metadata. Values outside the roadway are mapped to lane id 0 and are excluded from interaction mining. Algorithm~\ref{alg:lane-inference} summarises the procedure.

\begin{algorithm}[t]
\caption{Lane inference from lateral position and lane markings}
\label{alg:lane-inference}
\KwIn{Vehicle lateral position $y(t)$ and lane boundaries $\{b_0,\dots,b_K\}$ from recording metadata.}
\KwOut{Inferred lane index $\ell(t) \in \{1,\dots,K\}$ (or 0 if outside the roadway).}

\ForEach{frame $t$}{
  \If{$y(t) \notin [b_0, b_K)$}{
    $\ell(t) \leftarrow 0$\tcp*{outside valid lanes}
  }
  \Else{
    find $k$ such that $b_{k-1} \le y(t) < b_k$\;
    $\ell(t) \leftarrow k$\;
  }
}
\end{algorithm}

\subsection{Same-lane leader/follower reconstruction}
Within each frame and inferred lane, vehicles are ordered by longitudinal position in the travel direction. The immediate leader (ahead) and follower (behind) are given by adjacent vehicles in this ordering. Algorithm~\ref{alg:neighbor-reconstruction} summarises the procedure.

\begin{algorithm}[t]
\caption{Same-lane leader/follower reconstruction from geometry}
\label{alg:neighbor-reconstruction}
\KwIn{Frame-wise vehicle states with longitudinal position $x(t)$ and inferred lane index $\ell(t)$.}
\KwOut{For each vehicle: leaderId and followerId in the same lane (or 0 if none).}

\ForEach{frame $t$}{
  group vehicles by lane $\ell$\;
  \ForEach{lane group $\mathcal{V}_{\ell,t}$}{
    sort $\mathcal{V}_{\ell,t}$ by $x(t)$ in travel direction\;
    \ForEach{vehicle $v_i$ in sorted order}{
      leaderId$(v_i) \leftarrow$ closest vehicle ahead in the list (else 0)\;
      followerId$(v_i) \leftarrow$ closest vehicle behind in the list (else 0)\;
    }
  }
}
\end{algorithm}

These two components form the core of the \textbf{system under evaluation}: starting from trajectories and lane markings, we can reconstruct the relations needed for downstream interaction mining, indicator computation, and context-signature generation without using dataset-provided neighbour identifiers as inputs.


\section{Interaction mining: lane-change detection}
A lane change is accepted when a vehicle transitions from one stable lane block to another, where the source lane remains unchanged for at least $N_\text{before}$ frames before the transition and the target lane remains unchanged for at least $N_\text{after}$ frames after the transition. Transitions involving invalid lane ids (e.g., lane id 0) are ignored.

\section{Interaction mining: cut-in extraction (finding the follower)}
Given a lane change into target lane $L$ by vehicle $A$:
\begin{enumerate}[leftmargin=1.6em]
  \item starting from the lane change start, scan frames forward within a fixed search window,
  \item in each frame, identify a vehicle $B$ in lane $L$ for which $A$ is the immediate leader in that frame (according to the chosen source of neighbour relations: reconstructed relations or reference labels),
  \item if the relationship is established within the allowed delay and holds for at least a minimum number of frames, record a cut-in event with cutter $A$, follower $B$, and onset $t_0$ at the first matching frame.
\end{enumerate}

Algorithm~\ref{alg:cutin-detection} summarises the rule-based cut-in identification used in the reconstruction and reference-label paths.

\begin{algorithm}[t]
\caption{Cut-in identification as a cutter--follower pair}
\label{alg:cutin-detection}
\KwIn{Lane-change candidate for vehicle $A$ with target lane $L$, reconstructed relations over time.}
\KwOut{Cut-in event $(A,B,t_0)$ or no event.}

initialise $t_0 \leftarrow \text{undefined}$\;
\ForEach{frame $t$ in a search window after the lane change begins}{
  \If{$A$ is the leader of some vehicle $B$ in lane $L$}{
    set $t_0$ to the first such frame and keep $B$ as follower\;
    \If{the relation $A \rightarrow B$ persists for at least $d_{\min}$ frames}{
      \Return $(A,B,t_0)$\;
    }
  }
}
\Return no event\;
\end{algorithm}

\subsection{Design choices and their literature basis}
The framework keeps \emph{event extraction} separate from \emph{risk labelling}. The settings used in the final run are design choices, but they are informed by previous work rather than picked arbitrarily.

\begin{itemize}[leftmargin=1.2em]
  \item \textbf{Stable lane blocks.} We detect lane changes between stable lane segments rather than from single-frame label flips. This reduces sensitivity to label jitter and makes the onset definition reproducible.
  \item \textbf{Bounded pair formation.} A new follower must appear within a short delay and persist for a minimum duration. Naturalistic cut-in studies report wide variation in cut-in duration and gap acceptance \cite{wang2019cutin}, so persistence rules should be stated clearly instead of being hidden in later post-processing.
  \item \textbf{Pair-consistent indicator computation.} DHW, THW, and TTC are computed for the explicit cutter--follower pair, not for a generic dataset predecessor. This matters because recent lane-change risk models on \textit{highD} rely on interaction-specific features and time-to-lane-change rather than on anonymous vehicle states \cite{shangguan2022,liu2025multitask}.
  \item \textbf{Phase-aware interpretation.} Following Zhou and Itoh, TTC is especially informative when evaluating the targeted rear vehicle before or around insertion, while THW remains informative once the pair has been established in the target lane \cite{zhou2016riskperception}. For that reason, the thesis reports all three measures but uses execution-stage THW only as a transparent descriptive risk label.
\end{itemize}

These choices are not claimed to be universally best. Their purpose is to keep the cut-in definition clear, easy to inspect, and easy to vary in later sensitivity analysis.

\section{Safety quantification: surrogate measures and geometry model selection}
We compute DHW/THW/TTC for arbitrary pairs using geometry and speeds. Because trajectory datasets may define position references differently (e.g., centre point vs bounding-box corner), we \emph{select a geometry model} by comparing our computed indicators against highD-provided indicators for same-lane (ego, predecessor) pairs on a large random sample. We test multiple plausible reference conventions and select the one that best matches dataset-provided DHW/THW/TTC. In the final run, the best match was obtained with a bounding-box--based interpretation of the position variables (a top-left reference point), combined with vehicle length to compute front and rear positions.

To avoid negative gaps due to noise or reference mismatches, thesis-facing minima are clipped to be non-negative for DHW/THW. TTC is treated as infinite when the follower is not closing ($v_B \leq v_A$).

\section{Safety quantification: stage-wise features}
For each detected cut-in, we align a fixed time window around onset $t_0$ and compute stage-wise summaries for the four stages in Chapter~\ref{chap:theory}. The feature table includes stage minima of DHW/THW/TTC and kinematic aggregates for cutter and follower. These stage-wise summaries are later used not only for descriptive risk analysis, but also for operational-usefulness evaluation through threshold sensitivity and TTC-versus-THW disagreement analysis.

\section{Context signatures: encoding and reference verification}
To capture local traffic context around the cutter compactly, we encode neighbourhood occupancy into the 3$\times$3 semantic grid defined in Chapter~\ref{chap:theory}. Each frame is mapped to a binary occupancy pattern, and the 3$\times$3 cells are embedded into a 4$\times$4 grid and linearised using a Hilbert-type space-filling-curve order. To remove direction-induced asymmetry, events are mirror-normalised to a canonical orientation prior to aggregation.

\paragraph{Minimal-input SFC signatures.}
In the minimal-input path, occupancy is computed from inferred lanes (XY-lane) and lane-snapshot neighbour queries (geometry). This means SFC signatures can be generated without accessing highD neighbour columns.

\paragraph{Reference SFC signatures.}
For verification, we also construct a reference occupancy grid from the highD surrounding-vehicle identifiers (same-lane preceding/following plus left/right neighbours and alongside vehicles). Mapping each non-zero neighbour id to occupancy yields a reference 3$\times$3 grid per frame, which can be encoded into a reference SFC signature.

\paragraph{Signature agreement.}
We verify the SFC branch by comparing the reference and minimal-input codes for matching (event, frame) rows. We report both exact code match rate and the distribution of Hamming distances (bitwise differences), which localises whether mismatches affect a single cell or multiple cells.


Algorithm~\ref{alg:sfc-signature} summarises how the context signatures are constructed for each stage-aligned frame.

\begin{algorithm}[t]
\caption{Ego-centric occupancy grid and space-filling-curve context signature}
\label{alg:sfc-signature}
\KwIn{Cut-in event $(A,B,t_0)$, reconstructed lane/relations, neighbourhood range parameters.}
\KwOut{A stage-aligned sequence of binary grids and a 16-bit context signature per frame.}

\ForEach{frame $t$ in the analysis window around $t_0$}{
  build a 3$\times$3 ego-centric occupancy grid around the cutter $A$\;
  (optional) mirror-normalise the grid to a canonical direction\;
  embed the 3$\times$3 grid into a 4$\times$4 grid with fixed padding\;
  traverse cells in Hilbert order and write occupancy bits into a 16-bit code\;
  store the code and (optionally) the decoded grid for verification\;
}
\end{algorithm}

\section{Post-processing: reproducible outputs}
The framework is parameterised through a small set of explicit settings (e.g., stability thresholds and stage windows). Each run produces versioned outputs (tables and figures) that can be regenerated from the same raw data and settings. The intended entry point for reproducing the results is the project's reproducibility package or repository: \ArtifactURL.

% ==============================================================================
\chapter{Results}
\label{chap:results}

\section{Detected events and analysis coverage}
The final run uses recordings 01--60. Table~\ref{tab:events} summarises event yield.

\begin{table}[H]
\centering
\caption{Event counts for the final thesis run (recordings 01--60).}
\label{tab:events}
{\small
\setlength{\tabcolsep}{4pt}
\begin{tabularx}{\textwidth}{@{}XrrX@{}}
\toprule
Artefact & \shortstack{Lane\\changes} & \shortstack{Cut\\-ins} & Notes \\
\midrule
Interaction mining (stable blocks + follower pairing) & 11,856 & 10,026 & Lane changes and cut-ins extracted from reconstructed relations. \\
Stage feature table (intention--recovery window) & -- & 10,023 & Three events removed by stage-window validity checks. \\
Context signatures (frame-level occupancy) & -- & 1,421,980 & Frame-level context signatures for the 10,023 validated cut-in events. \\
\midrule
Recordings processed & \multicolumn{3}{c}{60 (01--60)} \\
\bottomrule
\end{tabularx}
}
\end{table}

Throughout the remainder of the thesis, we distinguish clearly between \emph{extracted cut-ins} (10,026) and \emph{validated analysis events} (10,023 after stage-window checks). This distinction is important because the framework exposes each acceptance criterion explicitly instead of silently dropping events in downstream scripts.

\section{Lane-change and cut-in duration statistics}
Table~\ref{tab:duration-stats} summarises dataset-wide duration statistics. Lane-change durations are computed from stable-block transitions. Cut-in relation duration is measured as the contiguous duration where the follower's predecessor is the cutter \emph{within the configured search window}. In our configuration, the relation duration is capped at 50 frames (\SI{2}{\second}), which explains the concentration at the cap.

\begin{table}[H]
\centering
\caption{Dataset-wide duration statistics (highD 01--60, \SI{25}{\hertz}).}
\label{tab:duration-stats}
{\small
\setlength{\tabcolsep}{4pt}
\begin{tabular}{@{}lrrrrrr@{}}
\toprule
Event type & Min (s) & P25 (s) & Median (s) & Mean (s) & P75 (s) & Max (s) \\
\midrule
Lane change duration & 0.960 & 3.680 & 6.240 & 6.551 & 8.960 & 69.120 \\
Cut-in relation duration & 0.360 & 1.960 & 1.960 & 1.870 & 1.960 & 1.960 \\
\bottomrule
\end{tabular}}
\end{table}

Overall, 84.56\% of detected lane changes are classified as cut-ins under our definition (10,026 cut-ins out of 11,856 lane changes).

\section{Evaluation of the minimal-input reconstruction (RQ1/RQ2)}
\paragraph{Same-lane neighbour reconstruction.}
Geometry-based same-lane neighbour reconstruction matches the reference preceding and following relations with mean accuracy of 0.999997 across the 60 recordings. Unless stated otherwise, accuracies are computed over all frames; conditional accuracies restricted to frames where the reference relation exists (non-zero id) are reported in the supplementary material.

\paragraph{XY-lane variant.}
When lane membership is also inferred from geometry (XY-lane), agreement remains high with mean preceding and following accuracies of 0.999810 and 0.999811. This corresponds to a frame-level mismatch rate of about 0.02\%. Cut-in detection based on reconstructed relations matches the reference-label events with $F_1 = 0.998938$.

These results answer the correctness question: the minimal-input path reproduces the reference-label path with high agreement. The remaining results address usefulness: how explicit pair construction, stage alignment, and explicit indicator choices affect interpretation once this correctness baseline is established.

\section{Execution-stage indicator distributions}
Table~\ref{tab:exec-dist} summarises execution-stage minima across all 10,023 events. TTC statistics are shown only for events with finite TTC (i.e., closing speed).

\begin{table}[H]
\centering
\caption{Execution-stage indicator minima across all events (10,023 cut-ins).}
\label{tab:exec-dist}
{\small
\setlength{\tabcolsep}{4pt}
\begin{tabular}{@{}lrrrrrr@{}}
\toprule
Indicator & Min & P25 & Median & Mean & P75 & Max \\
\midrule
DHW$_{\min}$ (m) & 1.37 & 20.69 & 32.21 & 51.28 & 59.82 & 373.98 \\
THW$_{\min}$ (s) & 0.051 & 0.769 & 1.194 & 1.700 & 2.028 & 12.383 \\
TTC$_{\min}$ (s), finite only (n=3,606) & 0.220 & 10.769 & 19.726 & 74.161 & 42.025 & 14428.50 \\
\bottomrule
\end{tabular}}
\end{table}

\section{Stage-based risk summaries}
Risk is defined as:
\[
\text{risky} := \min(\mathrm{THW})_{\text{execution}} < \SI{0.7}{\second}.
\]
Using this rule on the merged stage-feature table (10,023 events), 2,018 events are classified as risky (20.1\%). Table~\ref{tab:risk-summary} reports pooled comparison statistics. Here $\Delta v$ is defined per event as the difference between the cutter and follower mean speeds in the execution stage, i.e., $\Delta v = \bar v_{\text{cutter}}-\bar v_{\text{follower}}$ (reported as the median across events).

\begin{table}[H]
\centering
\caption{Execution-stage pooled risk summary (all recordings). $\Delta v$ is in \si{\meter\per\second}.}
\label{tab:risk-summary}
{\small
\setlength{\tabcolsep}{4pt}
\begin{tabular}{@{}lrrrr@{}}
\toprule
Group & Events & THW$_{\min}$ median (s) & DHW$_{\min}$ median (m) & $\Delta v$ median \\
\midrule
Non-risky & 8,005 & 1.427 & 39.680 & 2.304 \\
Risky & 2,018 & 0.539 & 14.945 & 2.853 \\
\bottomrule
\end{tabular}}
\end{table}

\section{Effect size comparison between risk groups}
Because the dataset contains over 10{,}000 cut-in events, even small group differences can appear statistically significant under conventional hypothesis tests. To describe \emph{magnitude} rather than only significance, Table~\ref{tab:effect-sizes} reports Cliff's delta ($\delta$) for selected execution-stage variables. Cliff's $\delta$ ranges from $-1$ (all risky values smaller than non-risky) to $+1$ (all risky values larger), and values near $0$ indicate substantial overlap between the groups.

\begin{table}[H]
\centering
\caption{Effect-size comparison between risk groups (execution stage). Medians are reported per group; $\delta$ is Cliff's delta (risky vs non-risky).}
\label{tab:effect-sizes}
\begin{tabular}{@{}lrrr@{}}
\toprule
Metric & Median non-risky & Median risky & Cliff's $\delta$ \\
\midrule
DHW$_{\min}$ (\si{\meter}) & 39.68 & 14.95 & $-0.950$ \\
$\Delta v$ (\si{\meter\per\second}) & 2.304 & 2.853 & $0.032$ \\
$\lvert v_{y,\text{cutter}}\rvert_{\max}$ (\si{\meter\per\second}) & 0.88 & 0.93 & $0.097$ \\
\bottomrule
\end{tabular}
\end{table}

The separation in DHW$_{\min}$ is expected because THW is directly related to spacing via $\mathrm{THW}=\mathrm{DHW}/v_B$. In contrast, the lateral-velocity difference is modest, indicating that the THW-based label primarily captures longitudinal gap acceptance rather than the lateral kinematics of the lane change.

\section{Sensitivity of the THW threshold}
To illustrate that the risk label is a design choice, Table~\ref{tab:sensitivity} shows how the risk prevalence changes with the THW threshold.

\begin{table}[H]
\centering
\caption{Sensitivity of risk prevalence to the THW execution-stage threshold (10,023 events).}
\label{tab:sensitivity}
\begin{tabular}{@{}lrr@{}}
\toprule
THW threshold (s) & Risky events & Risk ratio \\
\midrule
0.50 & 760 & 7.58\% \\
0.70 & 2,018 & 20.13\% \\
1.00 & 4,004 & 39.95\% \\
1.50 & 6,235 & 62.21\% \\
\bottomrule
\end{tabular}
\end{table}

\section{Indicator disagreement example: THW vs TTC}
TTC is only finite when the follower is closing in on the cutter ($v_B>v_A$). In this dataset, execution-stage TTC is finite for 3,606 of 10,023 events (35.98\%). Within this finite subset, risky events (by THW) have substantially lower TTC.

\begin{table}[H]
\centering
\caption{Execution-stage TTC for events with finite TTC only (3,606 events).}
\label{tab:ttc-finite}
\begin{tabular}{@{}lrrr@{}}
\toprule
Group (by THW label) & Events & TTC median (s) & TTC mean (s) \\
\midrule
Non-risky & 3,177 & 22.30 & 80.27 \\
Risky & 429 & 8.98 & 28.94 \\
\bottomrule
\end{tabular}
\end{table}

Table~\ref{tab:ttc-thw-overlap} compares the THW-based label to TTC thresholds over all events, treating infinite TTC as ``not risky''. Strict TTC thresholds have high precision but low recall relative to the THW label because many THW-risky events do not involve a closing-speed situation.

\begin{table}[H]
\centering
\caption{Overlap between THW-risky and TTC-threshold flags (all events; infinite TTC counted as not risky).}
\label{tab:ttc-thw-overlap}
\begin{tabular}{@{}rrrrrr@{}}
\toprule
TTC thr. (s) & TTC-flagged & TP & FP & Prec. & Rec. \\
\midrule
1.5 & 10 & 10 & 0 & 1.000 & 0.005 \\
3.0 & 43 & 37 & 6 & 0.860 & 0.018 \\
5.0 & 156 & 99 & 57 & 0.635 & 0.049 \\
10.0 & 788 & 237 & 551 & 0.301 & 0.117 \\
\bottomrule
\end{tabular}
\end{table}

\section{Why explicit pairing and phase-aware indicators matter}
The high agreement reported above answers the correctness question, but it does not by itself explain why the framework is useful. The extra value comes from making the event definition and the indicator interpretation explicit.

First, the event pipeline is traceable from end to end: 10,026 cut-ins were extracted, 10,023 remained after stage-window validity checks, and the stage feature table keeps the cutter and follower identities throughout the aligned window. This means the acceptance rules are visible in the method instead of being hidden inside dataset-specific scripts.

Second, Table~\ref{tab:sensitivity} and Table~\ref{tab:ttc-finite} show that one indicator cannot explain all phases equally well. When the execution-stage THW threshold is changed from 0.50 to 1.50 s, the risky-event share changes from 7.58\% to 62.21\%. At the same time, TTC is finite in only 35.98\% of execution-stage events. A TTC-only post-merge analysis would therefore ignore most established cut-ins even though the interacting pair is clearly defined.

This pattern is consistent with previous work showing that lane-change risk perception is stage dependent and that TTC and THW describe different parts of the interaction \cite{zhou2016riskperception,wang2021ssm}. In that sense, the contribution of \FrameworkName{} is not a claim that one threshold is universally correct. The contribution is that the framework reconstructs the right pair, computes the indicators for that pair, and keeps the later label choice explicit.

\section{Extreme events by TTC}
Table~\ref{tab:top-ttc} lists the three events with the lowest execution-stage TTC (finite only), which are useful for qualitative inspection using the dashboard in Figure~\ref{fig:gui-dashboard}.

\begin{table}[H]
\centering
\caption{Top-3 extreme events by minimum execution-stage TTC (finite only).}
\label{tab:top-ttc}
\begin{tabular}{@{}rrrrrrr@{}}
\toprule
Rec. & Cutter & Follower & From & To & TTC$_{\min}$ (s) & THW$_{\min}$ (s) \\
\midrule
40 & 682 & 687 & 1 & 2 & 0.220 & 0.051 \\
26 & 496 & 522 & 1 & 2 & 0.235 & 0.122 \\
30 & 839 & 852 & 1 & 2 & 0.288 & 0.067 \\
\bottomrule
\end{tabular}
\end{table}

\section{Per-recording variation}
Across the 60 recordings, the risk ratio (THW $<$ \SI{0.7}{\second}) varies between 11.1\% and 34.1\% (median 20.0\%). This suggests that traffic composition and local conditions differ across recordings even within the same dataset.

\section{Context signature maps}
\label{sec:context-maps}
To interpret the context signatures, we report mean occupancy probabilities in the semantic 3$\times$3 grid described in Chapter~\ref{chap:theory}. The grid is aligned with the cutter after mirror normalisation; columns correspond to \{left lane, same lane, right lane\} and rows to \{preceding, alongside, following\}. The centre cell (same-lane/alongside) is occupied by construction because it represents the cutter; occupancy in other cells indicates the presence of at least one surrounding vehicle in the corresponding relative region.

We present the decision stage as a representative example because it captures the traffic configuration immediately before the cut-in interaction becomes established. Figure~\ref{fig:sfc-decision-grids} shows mean occupancy for non-risk and risky events and the risk-minus-non-risk difference for the decision stage.

To highlight how group differences evolve over time, Figure~\ref{fig:sfc-stage-diff} shows risk-minus-non-risk difference maps for the intention, decision, and execution stages using a consistent colour scale. Because the differences are small in magnitude, they may appear visually subtle; supplementary absolute maps for the intention and execution stages are provided in Appendix~\ref{app:occupancy-supp}.

\begin{figure}[H]
\centering
\safeincludegraphics[width=0.31\textwidth]{decision_nonrisk.png}
\safeincludegraphics[width=0.31\textwidth]{decision_risk.png}
\safeincludegraphics[width=0.31\textwidth]{decision_diff.png}
\caption{Decision-stage occupancy maps: non-risk (left), risk (middle), and risk-minus-non-risk difference (right).}
\label{fig:sfc-decision-grids}
\end{figure}

\begin{figure}[H]
\centering
\safeincludegraphics[width=0.31\textwidth]{intention_diff.png}
\safeincludegraphics[width=0.31\textwidth]{decision_diff.png}
\safeincludegraphics[width=0.31\textwidth]{execution_diff.png}
\caption{Stage-wise risk-minus-non-risk occupancy differences for intention, decision, and execution stages.}
\label{fig:sfc-stage-diff}
\end{figure}

\section{Reference agreement for context signatures (RQ3)}
We verify the SFC encoding in two complementary ways:
\begin{itemize}[leftmargin=1.2em]
  \item \textbf{Algorithmic round-trip:} exhaustive decode$\leftrightarrow$encode verification across all $2^{16}$ codes confirms that the 16-bit SFC mapping is reversible for both Hilbert and Morton orders.
  \item \textbf{Reference agreement:} for cut-in stage frames, we compute a reference occupancy grid from highD surrounding-vehicle ids and compare its SFC code to the minimal-input code computed from inferred lanes and geometry. Exact matches indicate identical neighbour occupancy under the two definitions; non-matches localise differences to specific cells via Hamming distance.
\end{itemize}

Using the final thesis configuration (\SI{4.0}{\meter} alongside threshold; no fixed ahead/behind cutoff), we verify 1,421,980 stage rows and obtain 1,414,762 exact row matches (99.49\%). Stage-wise exact-match rates are 99.35\% (intention), 99.52\% (decision), and 99.60\% (execution). In the thesis event view, stage-window ANY-frame agreement is 100.00\% and stage-window ALL-frame agreement is 79.82\%. The raw unmirrored match rate is 57.56\%, which confirms that direction normalisation is required for like-for-like left/right occupancy comparison across opposite travel directions.

% ==============================================================================
\chapter{Discussion}
\label{chap:discussion}

\section{Answers to the research questions}
\textbf{RQ1 (reconstruction accuracy):} Same-lane neighbour reconstruction from geometry reaches near-perfect agreement with highD reference neighbours. With inferred lanes (XY-lane), the agreement remains above 0.9998 for both preceding and following.

\textbf{RQ2 (cut-in consistency):} Stable-block lane-change detection combined with follower pairing yields 10,026 extracted cut-ins from 11,856 lane changes across recordings 01--60, of which 10,023 remain after stage-window validity checks. When the full cut-in detector is run end-to-end on reconstructed relations, event matching against the reference-label path achieves $F_1 = 0.998938$. Beyond correctness, the explicit cutter--follower representation makes it possible to compute pair-consistent surrogate measures and to keep the event definition explicit.

\textbf{RQ3 (SFC encoding and verification):} A compact 3$\times$3 occupancy grid encoded into a 16-bit SFC signature provides an interpretable local context representation. The encoding is reversible by design and is validated by exhaustive round-trip checks. In the final configuration, reference-vs-minimal-input matching reaches 99.49\% exact row agreement over 1,421,980 verified rows (with 100\% event-level ANY-frame agreement), confirming that the reconstructed occupancy used for SFC encoding corresponds closely to the reference relations in highD when direction is normalised. Reference-vs-minimal-input signature matching provides an additional end-to-end verification that the reconstructed neighbour occupancy used for SFC encoding corresponds to the reference relations in highD.

\section{Interpretation and implications}
The main outcome of the thesis is methodological. What matters here is not the dashboard or the amount of code, but the fact that the method has a clear boundary. The framework uses trajectories and lane markings as inputs and treats highD relation fields as evaluation labels. Because of that, the same approach can be reused on data sources that do not provide neighbour ids, and disagreements with dataset conventions can be studied instead of inherited silently.

\begin{table}[H]
\centering
\caption{Main contribution claims and the evidence provided in this thesis.}
\label{tab:contribution-evidence}
{\small
\setlength{\tabcolsep}{4pt}
\begin{tabularx}{\textwidth}{@{}>{\raggedright\arraybackslash}p{0.24\textwidth} >{\raggedright\arraybackslash}p{0.28\textwidth} X@{}}
\toprule
Contribution claim & Evidence in this thesis & Why it matters \\
\midrule
Minimal-input relation reconstruction & Mean preceding/following accuracy 0.999810/0.999811; cut-in $F_1=0.998938$ & Makes cut-in mining reusable on data sources that lack neighbour ids. \\
Explicit cutter--follower pairing & 10,026 extracted cut-ins; 10,023 validated analysis events with stage-aligned pair features & Avoids reusing dataset indicators for the wrong interaction pair. \\
Phase-aware indicator analysis & TTC finite in 35.98\% of execution-stage events; THW sensitivity changes risk prevalence from 7.58\% to 62.21\% & Shows that risk interpretation depends on explicit pair and stage definitions rather than on one universal threshold. \\
Verified context signatures & Reversible 16-bit Hilbert code; 99.49\% exact row agreement over 1,421,980 stage rows & Provides compact, reversible context descriptors that can later support grouping or retrieval. \\
Clear evaluation boundary & Two-path design, parameter table, and deterministic exports & Separates method inputs from reference labels and makes the evaluation reproducible. \\
\bottomrule
\end{tabularx}}
\end{table}

Explicit pairing matters because the safety-relevant interaction after a cut-in is the cutter and the new follower in the target lane. Dataset-provided indicators computed for a generic predecessor relation may be correct for that dataset convention and still correspond to the wrong pair for a cut-in study. \FrameworkName{} addresses this by defining cut-ins as explicit cutter--follower events and computing DHW, THW, and TTC for that pair only.

The phase-aware indicator results matter for the same reason. TTC and THW are not interchangeable: TTC focuses on closing-speed conflict, while THW focuses on accepted gap. In our data, TTC is finite for only 35.98\% of execution-stage events, which means that a TTC-only post-insertion analysis would miss the majority of established cut-ins. This matches literature that treats lane-change risk perception as phase dependent \cite{zhou2016riskperception} and broader reviews that warn against one-size-fits-all interpretation of surrogate measures \cite{wang2021ssm}. The contribution here is therefore not a claim of one universally correct threshold, but a method that makes the indicator choice explicit and ties it to the correct interacting pair.

The SFC signatures complement the scalar indicators. Because the code is reversible, it does not replace interpretation with a black-box embedding. Instead, it provides a compact summary that can be decoded, checked, and later used for large-scale comparison, retrieval, or clustering of interaction patterns.

\section{Threats to validity}
\textbf{Construct validity:} Risk is defined via a THW threshold in the execution stage. Alternative definitions (e.g., TTC thresholds, required deceleration / braking-demand indicators, or post-encroachment measures) would label different subsets of events, so the reported ``risk ratio'' should be interpreted as \emph{prevalence under this definition}. More generally, our criteria should be read as a stricter, pair-consistent definition of cut-in relevance, not as proof that highD is absolutely ``wrong''. This matters especially because highD lacks wider upstream/downstream context that may influence discretionary lane changes \cite{wanggao2025highd}.

\textbf{Internal validity:} Reconstruction and metric computations depend on lane markings and coordinate conventions (e.g., bounding-box reference points). Errors in these inputs can shift onset $t_0$ and bias DHW/THW/TTC. We mitigate this via (i) geometry model selection against dataset-provided indicators where applicable, and (ii) independent reconstruction variants (geometry-only neighbours and XY-lane inference) that agree closely with the reference-label path.

\textbf{External validity:} Results are based on \textit{highD} highway recordings and may not transfer directly to urban driving, ramp merges, or different sensing setups. Differences in speed limits, lane geometry, and traffic composition can change both cut-in dynamics and the interpretation of surrogate measures.

\textbf{Conclusion validity:} We report descriptive differences, sensitivity checks, and agreement checks against reference labels. The thesis does not make causal claims about driver behaviour or safety outcomes, and statistical significance is not emphasised over effect magnitude.

\section{Future work}
Natural next steps include (i) extending the surrogate-measure layer with braking-demand or minimum-safe-deceleration criteria, which have been proposed as complementary lane-change safety indicators \cite{wang2020humanlike}, (ii) adding broader ablation studies over relevance-distance, pair-persistence, and stage-boundary criteria, ideally with manual review of disagreement cases, and (iii) exploiting the integer context signatures for larger-scale scenario mining, for example through retrieval, clustering, and cross-dataset comparison.

% ==============================================================================
\chapter{Conclusion}
\label{chap:conclusion}

This thesis presents \FrameworkName{}, a minimal-input framework for reconstructing lane and same-lane relations, detecting cut-ins as explicit cutter--follower events, computing pair-specific safety indicators, and encoding local traffic context as reversible SFC signatures. By using highD lane and neighbour ids only as reference labels, the thesis evaluates the method itself rather than simply analysing a dataset through its own derived fields.

On highD recordings 01--60, the framework reaches mean preceding/following accuracy 0.999810/0.999811 and event-level cut-in matching of $F_1 = 0.998938$. The detector extracts 10,026 cut-ins, of which 10,023 remain after stage-window checks. For 1,421,980 stage rows, the reconstructed and reference context signatures match exactly in 99.49\% of cases, and the 16-bit SFC encoding is verified by exhaustive encode$\leftrightarrow$decode round-trip checks.

The thesis also shows why this is more than a reconstruction exercise. Once the cutter and follower are made explicit, DHW, THW, and TTC can be computed for the correct pair and interpreted in a stage-based way. The threshold sensitivity analysis and TTC coverage results show that the later risk interpretation depends strongly on indicator choice and phase.

Overall, the contribution is a clear and reusable pipeline for cut-in analysis: minimal inputs, explicit interacting pair, verified context representation, and an evaluation setup that makes each step easy to inspect and reproduce. This gives a solid base for future retrieval, clustering, or cross-dataset use of cut-in scenarios.

% ==============================================================================
% Bibliography (self-contained, no BibTeX required)
% ==============================================================================
\cleardoublepage
\addcontentsline{toc}{chapter}{Bibliography}
\begin{thebibliography}{99}

\bibitem{highd_2018}
R.~Krajewski, J.~Bock, L.~Kloeker, and L.~Eckstein.
\newblock The highD dataset: A drone dataset of naturalistic vehicle trajectories on German highways for validation of highly automated driving systems.
\newblock In \emph{Proceedings of the IEEE 21st International Conference on Intelligent Transportation Systems (ITSC)}, 2018.
\newblock doi:10.1109/ITSC.2018.8569552.

\bibitem{hayward1972}
J.~C. Hayward.
\newblock Near-miss determination through use of a scale of danger.
\newblock \emph{Highway Research Record}, (384):24--34, 1972.

\bibitem{ssam2008}
D.~Gettman, L.~Pu, T.~Sayed, and S.~Shelby.
\newblock \emph{Surrogate Safety Assessment Model and Validation}.
\newblock Technical Report FHWA-HRT-08-051, Federal Highway Administration (FHWA), 2008.

\bibitem{minderhoud2001}
M.~M. Minderhoud and P.~H.~L. Bovy.
\newblock Extended time-to-collision measures for road traffic safety assessment.
\newblock \emph{Accident Analysis \& Prevention}, 33(1):89--97, 2001.

\bibitem{wang2021ssm}
C.~Wang, Y.~Xie, H.~Huang, and P.~Liu.
\newblock A review of surrogate safety measures and their applications in connected and automated vehicles safety modeling.
\newblock \emph{Accident Analysis \& Prevention}, 157:106157, 2021.
\newblock doi:10.1016/j.aap.2021.106157.

\bibitem{johnsson2021validation}
C.~Johnsson, A.~Laureshyn, and C.~D'Agostino.
\newblock A relative approach to the validation of surrogate measures of safety.
\newblock \emph{Accident Analysis \& Prevention}, 161:106350, 2021.
\newblock doi:10.1016/j.aap.2021.106350.

\bibitem{hwang2022fsmrl}
S.~Hwang, K.~Lee, H.~Jeon, and D.~Kum.
\newblock Autonomous vehicle cut-in algorithm for lane-merging scenarios via policy-based reinforcement learning nested within finite-state machine.
\newblock \emph{IEEE Transactions on Intelligent Transportation Systems}, 23(10):17594--17606, 2022.
\newblock doi:10.1109/TITS.2022.3153848.


\bibitem{wang2019cutin}
X.~Wang, M.~Yang, and D.~Hurwitz.
\newblock Analysis of cut-in behavior based on naturalistic driving data.
\newblock \emph{Accident Analysis \& Prevention}, 124:127--137, 2019.
\newblock doi:10.1016/j.aap.2019.01.006.

\bibitem{zhou2016riskperception}
H.~Zhou and M.~Itoh.
\newblock How does a driver perceive risk when making decision of lane-changing?
\newblock \emph{IFAC-PapersOnLine}, 49(19):60--65, 2016.
\newblock doi:10.1016/j.ifacol.2016.10.462.

\bibitem{shangguan2022}
Q.~Shangguan, T.~Fu, J.~Wang, S.~Fang, and L.~Fu.
\newblock A proactive lane-changing risk prediction framework considering driving intention recognition and different lane-changing patterns.
\newblock \emph{Accident Analysis \& Prevention}, 164:106500, 2022.
\newblock doi:10.1016/j.aap.2021.106500.

\bibitem{liu2025multitask}
Y.~Liu, J.~Zhang, N.~Lyu, and Q.~Zhao.
\newblock Predicting lane change maneuver and associated collision risks based on multi-task learning.
\newblock \emph{Accident Analysis \& Prevention}, 209:107830, 2025.
\newblock doi:10.1016/j.aap.2024.107830.

\bibitem{wang2020humanlike}
C.~Wang, Q.~Sun, Z.~Li, and H.~Zhang.
\newblock Human-like lane change decision model for autonomous vehicles that considers the risk perception of drivers in mixed traffic.
\newblock \emph{Sensors}, 20(8):2259, 2020.
\newblock doi:10.3390/s20082259.

\bibitem{kitchenham2007}
B.~Kitchenham and S.~Charters.
\newblock Guidelines for performing systematic literature reviews in software engineering.
\newblock Technical report, EBSE, 2007.

\bibitem{wohlin2014}
C.~Wohlin.
\newblock Guidelines for snowballing in systematic literature studies and a replication in software engineering.
\newblock In \emph{Proceedings of the 18th International Conference on Evaluation and Assessment in Software Engineering (EASE)}, 2014.

\bibitem{treiber2000idm}
M.~Treiber, A.~Hennecke, and D.~Helbing.
\newblock Congested traffic states in empirical observations and microscopic simulations.
\newblock \emph{Physical Review E}, 62(2):1805--1824, 2000.
\newblock (Introduces the Intelligent Driver Model, IDM.)

\bibitem{kesting2007mobil}
A.~Kesting, M.~Treiber, and D.~Helbing.
\newblock General lane-changing model MOBIL for car-following models.
\newblock \emph{Transportation Research Record}, 1999(1):86--94, 2007.

\bibitem{elfes1989}
A.~Elfes.
\newblock Using occupancy grids for mobile robot perception and navigation.
\newblock \emph{Computer}, 22(6):46--57, 1989.

\bibitem{deo2018}
N.~Deo and M.~M. Trivedi.
\newblock Convolutional social pooling for vehicle trajectory prediction.
\newblock In \emph{Proceedings of the IEEE/CVF Conference on Computer Vision and Pattern Recognition Workshops (CVPRW)}, 2018.

\bibitem{hilbert1891}
D.~Hilbert.
\newblock \"Uber die stetige Abbildung einer Linie auf ein Fl\"achenst\"uck.
\newblock \emph{Mathematische Annalen}, 38:459--460, 1891.

\bibitem{moon2001}
B.~Moon, H.~V. Jagadish, C.~Faloutsos, and J.~H. Saltz.
\newblock Analysis of the clustering properties of the Hilbert space-filling curve.
\newblock \emph{IEEE Transactions on Knowledge and Data Engineering}, 13(1):124--141, 2001.

\bibitem{wanggao2025highd}
Z.~Wang and Y.~Gao.
\newblock Lane change inconsistencies in the highD dataset.
\newblock \emph{Findings}, 2025.
\newblock doi:10.32866/001c.132341.

\end{thebibliography}

% ==============================================================================
% Appendix
% ==============================================================================
\appendix
\cleardoublepage

\chapter{Appendix: Terminology and highD field mapping}
\label{app:field-mapping}

To keep the main chapters focused on concepts rather than code-level details, we describe lanes and neighbour relations in descriptive terms such as \emph{same-lane leader} or \emph{right-lane follower}. Table~\ref{tab:field-mapping} lists the corresponding column names in the released highD dataset for readers who want to reproduce the evaluation directly from the raw files.

\begin{table}[H]
\centering
\caption{Mapping between thesis terminology and highD columns (tracks file).}
\label{tab:field-mapping}
{\small
\setlength{\tabcolsep}{4pt}
\begin{tabularx}{\textwidth}{@{}l l X@{}}
\toprule
Concept in this thesis & highD column & Meaning \\
\midrule
Lane index (reference labels) & \texttt{laneId} & Discrete lane identifier per frame. \\
Same-lane leader (ahead) & \texttt{precedingId} & Track id of closest vehicle ahead in the same lane (0 if none). \\
Same-lane follower (behind) & \texttt{followingId} & Track id of closest vehicle behind in the same lane (0 if none). \\
Left-lane leader & \texttt{leftPrecedingId} & Closest vehicle ahead in the left adjacent lane (0 if none). \\
Left-lane alongside & \texttt{leftAlongsideId} & Vehicle alongside in the left adjacent lane (0 if none). \\
Left-lane follower & \texttt{leftFollowingId} & Closest vehicle behind in the left adjacent lane (0 if none). \\
Right-lane leader & \texttt{rightPrecedingId} & Closest vehicle ahead in the right adjacent lane (0 if none). \\
Right-lane alongside & \texttt{rightAlongsideId} & Vehicle alongside in the right adjacent lane (0 if none). \\
Right-lane follower & \texttt{rightFollowingId} & Closest vehicle behind in the right adjacent lane (0 if none). \\
\bottomrule
\end{tabularx}}
\end{table}

HighD distributes each recording as three CSV files:
\begin{itemize}[leftmargin=1.2em]
  \item \texttt{XX\_tracks.csv} (per-frame trajectories),
  \item \texttt{XX\_tracksMeta.csv} (per-track metadata),
  \item \texttt{XX\_recordingMeta.csv} (recording metadata, including lane markings).
\end{itemize}

\chapter{Appendix: Plain-language algorithms}

\section{Lane change detection (stable blocks)}
\begin{enumerate}[leftmargin=1.6em]
  \item For each vehicle, read its lane id over time (reference) or inferred lane index (XY-lane).
  \item When the lane id changes from $L_1$ to $L_2$, check:
  \begin{itemize}[leftmargin=1.2em]
    \item the vehicle stayed in $L_1$ for at least $N_\text{before}$ frames before the change,
    \item the vehicle stayed in $L_2$ for at least $N_\text{after}$ frames after the change.
  \end{itemize}
  \item If both checks pass, emit a lane change event from $L_1$ to $L_2$.
\end{enumerate}

\section{Cut-in detection (how we find the follower)}
\begin{enumerate}[leftmargin=1.6em]
  \item Start from a lane change event where the cutter enters a target lane.
  \item Move forward frame by frame within a configured search window.
  \item In each frame, look for any vehicle in the target lane whose same-lane predecessor equals the cutter id.
  \item When this first happens, that vehicle is the follower and the frame is $t_0$.
  \item Require that the follower relation holds for a minimum number of frames to reduce noise.
\end{enumerate}

\section{DHW, THW, and TTC for the interacting pair}
\begin{itemize}[leftmargin=1.2em]
  \item DHW is computed as the longitudinal gap between the cutter's rear and the follower's front.
  \item THW is DHW divided by follower speed (with safe handling of near-zero speeds).
  \item TTC is DHW divided by closing speed when the follower is gaining on the cutter; otherwise TTC is treated as infinite.
\end{itemize}

\section{\texorpdfstring{Reference 3$\times$3}{Reference 3x3} occupancy grid from highD neighbour ids}
For a given ego vehicle and frame, highD provides surrounding-vehicle ids for left/right/preceding/following/alongside positions. We map these to the 3$\times$3 grid by setting a cell to $1$ if the corresponding neighbour id is non-zero; otherwise $0$. The same grid can be produced by geometry-based reconstruction, enabling code-level verification via SFC signature matching.



\chapter{Appendix: Supplementary occupancy maps}
\label{app:occupancy-supp}

Figure~\ref{fig:sfc-decision-grids} in Chapter~\ref{chap:results} presents the decision-stage occupancy maps in the main text because it provides a representative pre-onset view of the surrounding traffic configuration. For completeness, Figure~\ref{fig:sfc-supp-abs} reports additional absolute occupancy maps for the intention and execution stages.

\begin{figure}[H]
\centering
\safeincludegraphics[width=0.48\textwidth]{intention_nonrisk.png}
\safeincludegraphics[width=0.48\textwidth]{intention_risk.png}\\[2pt]
\safeincludegraphics[width=0.48\textwidth]{execution_nonrisk.png}
\safeincludegraphics[width=0.48\textwidth]{execution_risk.png}
\caption{Supplementary absolute occupancy maps for the intention and execution stages (non-risk vs risky). These plots complement the decision-stage example in Figure~\ref{fig:sfc-decision-grids}.}
\label{fig:sfc-supp-abs}
\end{figure}


\end{document}
